% lem-ballpotential -- CERW.Frozen.norm_ball_potential (draft). Source: limit-shapes.tex:361-366, label lem:ballpotential
\paragraph{Paper (verbatim).}
\begin{lemma}[The potential of a norm ball]\label{lem:ballpotential}
Let $\gauge$ be a norm, let $\rho>0$, and let $y\in\R^d$. Then
\begin{equation}\label{eq:ballpotential}
U_{\{\gauge<\rho\}}(y)=2d\drift(\rho-\gauge(y))_+\, .
\end{equation}
\end{lemma}
Definitions the statement relies on (verbatim):
\textit{eq:potential-norm (lines 349--353)}
For a norm~$\gauge$, the potential~\eqref{eq:potential-intro} is defined with the gradient of~$\gauge$ in place of~$v/|v|$: for a bounded measurable set~$D\subset\R^d$ and $y\in\R^d$, let
\begin{equation}\label{eq:potential-norm}
U_D(y)\coloneqq\frac{2\drift}{\omega_d}\int_D\nabla\gauge(v)\cdot\frac{v-y}{|v-y|^d}\dd v
\end{equation}
be the potential generated by the vector field~$-\drift\nabla\gauge$ on~$D$; for the Euclidean norm it is~\eqref{eq:potential-intro}.
\paragraph{Transcription.}
Paper macros: \verb!\drift! is $\kappa$ (Lean \verb!ε!), \verb!\gauge! is $\Psi$.

\emph{Object mapping.}
\begin{itemize}
\item $\gauge$: \verb!Ψ! with \verb!hΨ : CERW.IsNorm Ψ!; $\rho>0$: \verb!{ρ : ℝ} (hρ : 0 < ρ)!; $y\in\R^d$: \verb!y : EuclideanSpace ℝ (Fin d)!; $d\ge2$ (standing): \verb!hd : 2 ≤ d!.
\item $U_D$ of eq:potential-norm is \verb!CERW.normPotential d ε Ψ D y :=  2 * ε / unitBallVolume d * ∫ v in D, inner ℝ (gradient Ψ v) (v - y) / ‖v - y‖ ^ d!, with \verb!unitBallVolume d = (volume (Metric.ball 0 1)).toReal! $=\omega_d$ and \verb!‖v - y‖ ^ d! the natural-number power $|v-y|^d$; the integrand is $\nabla\gauge(v)\cdot(v-y)/|v-y|^d$.
\item $\nabla\gauge$ is Mathlib's \verb!gradient Ψ v!, the gradient where $\gauge$ is differentiable and $0$ elsewhere. A norm is Lipschitz, so the non-differentiability set is Lebesgue-null (Rademacher; machine-checked \verb!ae_differentiable!), and \verb!gradient Ψ! occurs only under an integral, so the value $0$ is never seen. Moreover \verb!gradient Ψ! is Borel measurable and bounded by the Lipschitz constant $K$, the integrand is bounded by $K|v-y|^{1-d}$, and this is integrable on every bounded measurable set (machine-checked \verb!integrableOn_normPotential_integrand!): the Bochner integral is a genuine integral. At $v=y$ the integrand reads $0/0=0$, a single null point.
\item The ball $\{\gauge<\rho\}$ is \verb!{v | Ψ v < ρ}! (open and bounded). $(\rho-\gauge(y))_+$ is \verb!max (ρ - Ψ y) 0! and $2d\kappa$ is \verb!2 * d * ε!.
\end{itemize}

\emph{Reading choices.}
\begin{enumerate}
\item \textbf{$\kappa$ is an arbitrary real.} In the paper $\kappa$ satisfies the standing ellipticity condition, but the lemma does not use it. The draft binds \verb!(ε : ℝ)! with no hypothesis. This is true for every real $\varepsilon$, because \verb!normPotential d ε Ψ D y = ε * normPotential d 1 Ψ D y! and the right side is $\varepsilon\cdot2d(\rho-\gauge(y))_+$ (machine-checked: \verb!ball_all_eps! derives all $\varepsilon$ from $\varepsilon=1$). It is therefore a true strengthening of the paper's lemma; consumers at an admissible $\varepsilon$ instantiate it.
\item The identity is an equality of real numbers; no \verb!ℝ≥0∞! is involved.
\item Euclidean consistency (machine-checked): for $\gauge=|\cdot|$ one has \verb!gradient ‖·‖ v = unitDir v! at every $v$ (value $0$ at $v=0$, as $u_0=0$), hence \verb!normPotential d ε ‖·‖ D y = potential d ε D y!; and the statement at $\varepsilon>0$ is the ball clause of the proved \verb!CERW.Frozen.potential_geometry!.
\item Numerical check of the normalisation ($d=2$, $\gauge=\|\cdot\|_\infty$, $\rho=1$, $\varepsilon=1$): $U(0,0)\approx3.9988$ against $4$, and $U(0.3,0.2)\approx2.798$ against $2.8$.
\end{enumerate}
